\documentclass[12pt,a4paper]{article}
\usepackage[utf8]{inputenc}
\usepackage[indonesian]{babel}
\usepackage{geometry}
\geometry{a4paper, margin=3cm}
\usepackage{graphicx}
\usepackage{titlesec}
\usepackage{setspace}
\usepackage{listings}
\usepackage{xcolor}

% Setup untuk style kode
\definecolor{codegreen}{rgb}{0,0.6,0}
\definecolor{codegray}{rgb}{0.5,0.5,0.5}
\definecolor{codepurple}{rgb}{0.58,0,0.82}
\definecolor{backcolour}{rgb}{0.95,0.95,0.92}

\lstdefinestyle{mystyle}{
    backgroundcolor=\color{backcolour},   
    commentstyle=\color{codegreen},
    keywordstyle=\color{magenta},
    numberstyle=\tiny\color{codegray},
    stringstyle=\color{codepurple},
    basicstyle=\ttfamily\footnotesize,
    breakatwhitespace=false,         
    breaklines=true,                 
    captionpos=b,                    
    keepspaces=true,                 
    numbers=left,                    
    numbersep=5pt,                  
    showspaces=false,                
    showstringspaces=false,
    showtabs=false,                  
    tabsize=2
}
\lstset{style=mystyle}

% Definisi bahasa JSON sederhana untuk listing
\lstdefinelanguage{json}{
    basicstyle=\normalfont\ttfamily\footnotesize,
    numbers=left,
    numberstyle=\tiny\color{codegray},
    stepnumber=1,
    numbersep=5pt,
    showstringspaces=false,
    breaklines=true,
    string=[s]{"}{"},
    stringstyle=\color{codepurple},
    comment=[l]{:},
    commentstyle=\color{black},
}

\begin{document}

% ------------------ COVER PAGE ------------------
\begin{titlepage}
    \centering
    \vspace*{2cm}
    
    {\Large \textbf{Source Code Hak Cipta}\par}
    \vspace{0.5cm}
    {\large \textbf{(PROGRAM KOMPUTER)}\par}
    
    \vspace{2.5cm}
    
    {\Large \textbf{Golden Hour: Sistem Game Simulasi Edukasi Pertolongan Pertama Pasien Trauma Fisik Berbasis Keputusan dan Aksi}\par}
    
    \vspace{3cm}
    
    {\large \textbf{TIM PENYUSUN}\par}
    \vspace{0.5cm}
    \begin{center}
    1. Achmad Arifin S.T., M.Eng., Ph.D. \\
    2. Rasyida Shabihah Zukro Aini, S.T., M.T. \\
    3. Dr. Norma Hermawan, S.T., M.T., M.Sc. \\
    4. Rizki Tetuko Irianto \\
    \end{center}
    
    \vfill
    
    {\large \textbf{INSTITUT TEKNOLOGI SEPULUH NOPEMBER SURABAYA}\par}
    {\large \textbf{2026}\par}
    
    \vspace*{1cm}
\end{titlepage}

% ------------------ CONTENT ------------------
\newpage
\tableofcontents
\newpage

\section{Deskripsi Sistem}
Sistem perangkat lunak ini merupakan bagian dari aplikasi game simulasi edukasi medis "Golden Hour". Kode sumber utama yang menyusun Hak Cipta Program Komputer ini terdiri dari arsitektur \textit{Game Manager} yang mengontrol aliran logika permainan, \textit{OllamaManager} untuk pemrosesan interaksi cerdas, \textit{QuizData} sebagai struktur cetak biru skenario, dan file konfigurasi JSON sebagai pangkalan data cerita dan aksi medis interaktif.

\section{Arsitektur Inti: GameManager.cs}
Modul ini adalah otak dari seluruh simulasi. Mengelola pergantian skenario, validasi kuis, mekanisme \textit{cutscene}, dan transisi UI interaktif.
\lstinputlisting[language={[Sharp]C}]{Assets/Scripts/GameManager.cs}

\newpage
\section{Struktur Data: QuizData.cs}
Modul ini mendefinisikan kelas serialisasi yang menjembatani data skenario (JSON) ke dalam objek C\# yang diproses oleh \textit{engine} Unity.
\lstinputlisting[language={[Sharp]C}]{Assets/Scripts/QuizData.cs}

\newpage
\section{Integrasi AI: OllamaManager.cs}
Modul inovatif yang menangani komunikasi HTTP asinkron ke server model bahasa lokal (Ollama) untuk menghasilkan percakapan atau sistem \textit{feedback} yang dinamis di dalam simulasi.
\lstinputlisting[language={[Sharp]C}]{Assets/Scripts/OllamaManager.cs}

\newpage
\section{Basis Data Skenario Medis: Scenario1.json}
File struktur data yang menyimpan urutan instruksi medis, paramater kuis (seperti ambang batas waktu dan tingkat fatalitas kesalahan), serta penentu aset visual untuk setiap transisi simulasi.
\lstinputlisting[language=json]{Assets/Resources/Scenario1.json}

\end{document}
